\documentclass[12pt]{article}
\usepackage[utf8]{inputenc}
\usepackage[spanish]{babel}
\usepackage[a4paper, margin=2cm]{geometry}
\usepackage{tikz}
\usetikzlibrary{babel}
\usepackage[most]{tcolorbox}
\usepackage{amsmath}

% Usar fuente sin serifa para un aspecto más moderno (tipo infografía)
\renewcommand{\familydefault}{\sfdefault}

% Definición de colores con temática médica/científica
\definecolor{medblue}{RGB}{43, 108, 176}
\definecolor{medteal}{RGB}{49, 151, 149}
\definecolor{medlight}{RGB}{235, 248, 255}
\definecolor{meddark}{RGB}{26, 32, 44}
\definecolor{pillred}{RGB}{239, 68, 68}    % Rojo para la cápsula
\definecolor{scalebg}{RGB}{209, 213, 219}  % Gris para la báscula

\begin{document}
\pagestyle{empty}

\begin{tcolorbox}[
    enhanced, 
    colback=medlight!30, 
    colframe=medblue, 
    title={\Large \textbf{Sistemas de Unidades: Dosis Médica por Peso}}, 
    halign title=center, 
    fonttitle=\bfseries, 
    arc=4mm, 
    boxrule=1.5pt,
    shadow={2mm}{-2mm}{0mm}{black!30!white}
]

    \vspace{0.2cm}
    \begin{tcolorbox}[colback=white, colframe=medteal, arc=2mm, boxrule=1.2pt, title=\textbf{Caso Clínico / Problema}]
        Un paciente recibe un medicamento a razón de \textbf{5 mg por cada kg} de peso corporal. Si en el expediente clínico el paciente registra un peso de \textbf{154 libras (lb)}. Sabiendo que \textbf{1 kg $\approx$ 2.2 lb}, ¿cuántos miligramos del medicamento debe recibir?
    \end{tcolorbox}

    \vspace{0.4cm}
    
    \begin{center}
    \begin{tikzpicture}[scale=0.95]
        
        % --- Paciente en Libras (Izquierda) ---
        \begin{scope}[shift={(-4, 0)}]
            % Báscula
            \filldraw[scalebg, draw=meddark, thick, rounded corners=2pt] (-1.2,-1.2) rectangle (1.2,-0.8);
            \node[font=\bfseries, text=meddark, scale=0.9] at (0,-1) {154 lb};
            % Cuerpo del paciente estilizado
            \filldraw[medblue!50, draw=meddark, thick] (0, 1.2) circle (0.5cm);
            \filldraw[medblue!50, draw=meddark, thick, rounded corners=4pt] (-0.8, -0.8) rectangle (0.8, 0.6);
            
            \node[font=\bfseries, text=meddark, align=center] at (0, -2) {Peso Inicial\\(Libras)};
        \end{scope}
        
        % --- Conversión 1 ---
        \draw[->, ultra thick, medteal] (-2.5, 0.2) -- (-0.5, 0.2) node[midway, above, font=\small\bfseries] {$\div 2.2$};
        
        % --- Paciente en Kilogramos (Centro) ---
        \begin{scope}[shift={(1, 0)}]
            % Báscula
            \filldraw[medteal!20, draw=medteal, thick, rounded corners=2pt] (-1.2,-1.2) rectangle (1.2,-0.8);
            \node[font=\bfseries, text=medteal!80!black, scale=0.9] at (0,-1) {70 kg};
            % Cuerpo del paciente estilizado
            \filldraw[medteal!50, draw=medteal!80!black, thick] (0, 1.2) circle (0.5cm);
            \filldraw[medteal!50, draw=medteal!80!black, thick, rounded corners=4pt] (-0.8, -0.8) rectangle (0.8, 0.6);
            
            \node[font=\bfseries, text=meddark, align=center] at (0, -2) {Peso Clínico\\(Kilogramos)};
        \end{scope}
        
        % --- Conversión 2 ---
        \draw[->, ultra thick, pillred] (2.5, 0.2) -- (4.5, 0.2) node[midway, above, font=\small\bfseries] {$\times 5 \text{ mg/kg}$};
        
        % --- Medicamento / Dosis (Derecha) ---
        \begin{scope}[shift={(6, 0)}]
            % Cápsula rotada para dinamismo
            \begin{scope}[rotate=30]
                \filldraw[pillred, draw=meddark, thick] (-0.9, -0.4) -- (0, -0.4) -- (0, 0.4) -- (-0.9, 0.4) arc (90:270:0.4cm);
                \filldraw[white, draw=meddark, thick] (0, -0.4) -- (0.9, -0.4) arc (-90:90:0.4cm) -- (0, 0.4) -- cycle;
                % Detalle de brillo en la cápsula
                \fill[white, opacity=0.6] (-0.7, 0.1) rectangle (-0.1, 0.25);
            \end{scope}
            
            \node[font=\Large\bfseries, text=pillred!80!black, align=center] at (0, 1.5) {350 mg};
            \node[font=\bfseries, text=meddark, align=center] at (0, -2) {Dosis Total\\Terapéutica};
        \end{scope}
        
    \end{tikzpicture}
    \end{center}

    \vspace{0.4cm}
    \begin{tcolorbox}[colback=medteal!10, colframe=medteal, arc=2mm, boxrule=1.2pt, title=\textbf{Desarrollo Matemático y Solución}]
        Para calcular la dosis correcta, el proceso se divide en dos etapas:
        
        \textbf{Paso 1: Convertir el peso del paciente a kilogramos (kg).}
        \[ \text{Peso} = 154 \text{ lb} \times \left( \frac{1 \text{ kg}}{2.2 \text{ lb}} \right) = \mathbf{70 \text{ kg}} \]
        
        \textbf{Paso 2: Calcular la dosis terapéutica en base al peso.}
        \[ \text{Dosis} = 70 \text{ kg} \times \left( \frac{5 \text{ mg}}{1 \text{ kg}} \right) = \mathbf{350 \text{ mg}} \]
        
        \textbf{Resultado:} El paciente debe recibir exactamente \textbf{350 mg} del medicamento.
    \end{tcolorbox}
    
\end{tcolorbox}

\end{document}